% English translation, made from our own transcription (original.tex), not from any existing
% translation. Every math expression, symbol, \tag{}, \origpage{N} marker, and \ednote is
% reproduced unchanged; only the prose is translated — including the short prose set INSIDE math
% via text inserts (HOUSESTYLE R21):
%   * \text{durch} → \text{by} (p. 161)
%   * array condition labels in display (1.) on p. 165
%   * definition clauses in display (1.) on p. 166
%   * ordinal suffixes: "sten" / "ten" / "ter" → "th" ($54^{\text{th}}$, $n^{\text{th}}$, $m^{\text{th}}$,
%     $(m+1)^{\text{th}}$).
% pipeline/texcompare.py ignores the CONTENT of such \text inserts while checking that each insert
% is present, unmoved, and that all surrounding math is identical.
%
% TERMINOLOGY. Follows corpus/riemann-1866-verschwinden-theta-functionen/glossary.yaml and
% corpus/riemann-1857-abelsche-functionen/glossary.yaml:
%   * Abtheilung = "Division"
%   * Abhandlung = "memoir"
%   * Querschnitt = "crosscut"
%   * gleichverzweigt = "equally branched"
%   * Verzweigungsart = "mode of branching"
%   * Modulsystem = "system of moduli"
%   * unendlich klein von der ersten/zweiten Ordnung = "infinitely small of the first/second order"
%   * Grössensystem = "system of quantities", Grössengebiet = "domain of quantities"
%   * Werthsystem = "system of values"
%   * Bestimmungsweise = "mode of determination"
%   * identisch verschwinden = "vanish identically"
%
% The printer's misprint flagged with an \ednote on p. 171 (a product sign \Pi in the denominator
% of formula (4.) where a summation \Sigma over \nu is required) is carried over verbatim as printed.
\documentclass{article}
\usepackage{readmasters}
\begin{document}

\origpage{161}
\section*{On the Vanishing of the $\vartheta$-Functions.}

(By Mr.~\emph{B.~Riemann} in Göttingen.)

The second Division of my theory of \emph{Abel}ian functions, published in the $54^{\text{th}}$ volume of this Journal, contains the proof of a theorem on the vanishing of the $\vartheta$-functions, which I shall immediately restate, assuming the notations employed in that memoir to be known to the reader. Everything that still follows in the memoir contains brief indications concerning the application of this theorem, which, under our method—based upon the determination of functions by their discontinuities and their becoming infinite—must, as is readily seen, form the foundation of the theory of \emph{Abel}ian functions. In the theorem itself and its proof, however, the circumstance was not properly taken into account that the $\vartheta$-function can vanish identically, i.e.\ for every value of this variable, through the substitution of the integrals of algebraic functions of a single variable. To remedy this deficiency is the purpose of the following short memoir.

In the presentation of investigations on $\vartheta$-functions with an indeterminate number of variables, the need for an abbreviated notation for a sequence, such as
\[
v_1,\quad v_2,\quad \dots,\quad v_m
\]
makes itself felt as soon as the expression of $v_\nu$ in terms of $\nu$ is complicated. One could form this symbol in complete analogy with the summation and product signs; such a notation would, however, take up too much space and within the functional symbols be inconvenient for printing; I therefore prefer to denote
\[
v_1,\quad v_2,\quad \dots,\quad v_m \quad\text{by}\quad \left(\overset{m}{\underset{1}{v}}(v_\nu)\right)
\]
and thus
\[
\vartheta(v_1, v_2, \dots, v_p) \quad\text{by}\quad \vartheta\left(\overset{p}{\underset{1}{v}}(v_\nu)\right).
\]

\section*{1.}

If in the function $\vartheta(v_1, v_2, \dots, v_p)$ one substitutes for the $p$ variables $v$ the $p$ integrals $u_1-e_1$, $u_2-e_2$, \dots, $u_p-e_p$ of algebraic functions of $z$ branched like the surface
\origpage{162}
$T$, one obtains a function of $z$ which varies continuously throughout the surface $T$ except across the lines $b$, but upon passing from the negative to the positive side of the line $b_\nu$ acquires the factor $e^{-u^+_\nu-u^-_\nu+2e_\nu}$. As was proved in \S.~22, this function, if it does not vanish for all values of $z$, becomes infinitely small of the first order for only $p$ points of the surface $T$. These points were denoted by $\eta_1$, $\eta_2$, \dots, $\eta_p$, and the value of the function $u_\nu$ at the point $\eta_\mu$ by $\alpha_\nu^{(\mu)}$. There then resulted, with respect to the $2p$ systems of moduli of the $\vartheta$-function, the congruence
\[
\tag{(1.)}
(e_1, e_2, \dots, e_p) \equiv \left(\sum_1^p \alpha_1^{(\mu)} + K_1,\, \sum_1^p \alpha_2^{(\mu)} + K_2,\, \dots,\, \sum_1^p \alpha_p^{(\mu)} + K_p\right),
\]
wherein the quantities $K$ depended on the hitherto still arbitrary additive constants in the functions $u$, but were independent of the quantities $e$ and the points $\eta$.

If one carries out the calculation indicated there, one finds
\[
\tag{(2.)}
2K_\nu = \sum \frac{1}{\pi i} \int (u_\nu^+ + u_\nu^-) \, du_{\nu'} - \varepsilon_\nu \pi i - \sum_{\mu=1}^{\mu=p} \varepsilon'_\mu a_{\mu,\nu}.
\]
In this expression the integral $\displaystyle\int (u_\nu^+ + u_\nu^-) \, du_{\nu'}$ is to be extended positively along $b_{\nu'}$, and in the sum all numbers from 1 to $p$ except $\nu$ are to be taken for $\nu'$; $\varepsilon_\nu = \pm 1$, according as the end of $l_\nu$ lies on the positive or negative side of $a_\nu$, and $\varepsilon'_\nu = \pm 1$, according as the same lies on the positive or negative side of $b_\nu$. The determination of the signs is, moreover, only necessary if the quantities $e$ are to be completely determined from the discontinuities of $\log \vartheta$ according to the equations given in \S.~22; the above congruence (1.) remains valid whatever signs one may choose.

We retain for the present the simplifying assumption made there, that the additive constants in the functions $u$ are determined so that the quantities $K$ are all equal to zero. In order to free the results thus obtained ultimately from this restrictive assumption, one obviously needs only to add $-K_1$, $-K_2$, \dots, $-K_p$ to the arguments everywhere in the $\vartheta$-functions.

If therefore the function $\vartheta(u_1-e_1, u_2-e_2, \dots, u_p-e_p)$ vanishes for the $p$ points $\eta_1$, $\eta_2$, \dots, $\eta_p$ \emph{and does not vanish identically for every value of $z$}, then
\[
(e_1, e_2, \dots, e_p) \equiv \left(\sum_1^p \alpha_1^{(\mu)},\, \sum_1^p \alpha_2^{(\mu)},\, \dots,\, \sum_1^p \alpha_p^{(\mu)}\right).
\]

\origpage{163}
This theorem holds for entirely arbitrary values of the quantities $e$, and from this, by letting the point $(s, z)$ coincide with the point $\eta_p$, we concluded that
\[
\vartheta\left(-\sum_1^{p-1} \alpha_1^{(\mu)}, -\sum_1^{p-1} \alpha_2^{(\mu)}, \dots, -\sum_1^{p-1} \alpha_p^{(\mu)}\right) = 0,
\]
or, since the $\vartheta$-function is even,
\[
\vartheta\left(\sum_1^{p-1} \alpha_1^{(\mu)}, \sum_1^{p-1} \alpha_2^{(\mu)}, \dots, \sum_1^{p-1} \alpha_p^{(\mu)}\right) = 0,
\]
whatever the points $\eta_1$, $\eta_2$, \dots, $\eta_{p-1}$ may be.

\section*{2.}

The proof of this theorem requires, however, a completion because of the circumstance that the function $\vartheta(u_1-e_1, u_2-e_2, \dots, u_p-e_p)$ can vanish identically (which indeed occurs for certain values of the quantities $e$ in every system of equally branched algebraic functions).

Because of this circumstance one must be content, first, to show that the theorem remains valid while the points $\eta$ change their positions independently of one another within finite limits. From this the general validity of the theorem then follows according to the principle that a function of a complex quantity cannot be equal to zero within a finite domain without being equal to zero everywhere.

If $z$ is given, the quantities $e_1$, $e_2$, \dots, $e_p$ can always be chosen so that $\vartheta(u_1-e_1, u_2-e_2, \dots, u_p-e_p)$ does not vanish; for otherwise the function $\vartheta(v_1, v_2, \dots, v_p)$ would have to vanish for all values of the quantities $v$, and consequently in its expansion in integral powers of $e^{2v_1}$, $e^{2v_2}$, \dots, $e^{2v_p}$ all coefficients would have to be equal to zero, which is not the case. The quantities $e$ can then vary independently of one another within finite domains of quantities without the function $\vartheta(u_1-e_1, u_2-e_2, \dots, u_p-e_p)$ vanishing for this value of $z$. Or in other words: one can always specify a domain of quantities $E$ of $2p$ dimensions within which the system of quantities $e$ can move without the function $\vartheta(u_1-e_1, u_2-e_2, \dots, u_p-e_p)$ vanishing for this value of $z$. It therefore becomes infinitely small of the first order for only $p$ positions of $(s, z)$, and if one denotes these points by $\eta_1$, $\eta_2$, \dots, $\eta_p$, then
\[
\tag{(1.)}
(e_1, e_2, \dots, e_p) \equiv \left(\sum_1^p \alpha_1^{(\mu)}, \sum_1^p \alpha_2^{(\mu)}, \dots, \sum_1^p \alpha_p^{(\mu)}\right).
\]

\origpage{164}
To every mode of determination of the system of quantities $e$ within $E$, or to every point of $E$, there then corresponds a mode of determination of the points $\eta$, the totality of which forms a domain of quantities $H$ corresponding to the domain of quantities $E$. In consequence of equation (1.), however, to each point of $H$ there also corresponds only one point of $E$; if then $H$ had only $2p-1$, or fewer dimensions, $E$ could not have $2p$ dimensions. Consequently $H$ has $2p$ dimensions. The inferences on which our theorem rests therefore remain applicable for arbitrary positions of the points $\eta$ within finite domains, and the equation
\[
\vartheta\left(-\sum_1^{p-1} \alpha_1^{(\mu)},\, -\sum_1^{p-1} \alpha_2^{(\mu)},\, \dots,\, -\sum_1^{p-1} \alpha_p^{(\mu)}\right) = 0
\]
holds for arbitrary positions of the points $\eta_1$, $\eta_2$, \dots, $\eta_{p-1}$ within finite domains, and consequently in general.

\section*{3.}

From this it follows that the system of quantities $(e_1, e_2, \dots, e_p)$ can always and only in one way be set congruent to an expression of the form $\left(\overset{p}{\underset{1}{v}}\left(\displaystyle\sum_1^p \alpha_\nu^{(\mu)}\right)\right)$ when $\vartheta\left(\overset{p}{\underset{1}{v}}(u_\nu - e_\nu)\right)$ does not vanish for every value of $z$; for if the points $\eta_1$, $\eta_2$, \dots, $\eta_p$ could be determined in more than one way so that the congruence
\[
\left(\overset{p}{\underset{1}{v}}(e_\nu)\right) \equiv \left(\overset{p}{\underset{1}{v}}\left(\sum_1^p \alpha_\nu^{(\mu)}\right)\right)
\]
were satisfied, then according to the theorem just proved the function $\vartheta\left(\overset{p}{\underset{1}{v}}(u_\nu - e_\nu)\right)$ would vanish for more than $p$ points without being identically equal to zero, which is impossible.

If $\vartheta\left(\overset{p}{\underset{1}{v}}(u_\nu - e_\nu)\right)$ vanishes identically, one must, in order to put $\left(\overset{p}{\underset{1}{v}}(e_\nu)\right)$ into the above form, consider $\vartheta\left(\overset{p}{\underset{1}{v}}(u_\nu + \alpha_\nu^{(1)} - u_\nu^{(1)} - e_\nu)\right)$, and if this function vanishes identically for every value of $z$, $s_1$, $z_1$, the function $\vartheta\left(\overset{p}{\underset{1}{v}}\left(u_\nu + \displaystyle\sum_1^2 \alpha_\nu^{(\mu)} - \displaystyle\sum_1^2 u_\nu^{(\mu)} - e_\nu\right)\right)$.

\origpage{165}
\[
\tag{(1.)}
\left\{
\begin{array}{l}
\text{We assume that} \\
\quad \vartheta\left(\overset{p}{\underset{1}{v}}\left(\sum_1^m \alpha_\nu^{(p+2-\mu)} - \sum_1^{m-1} u_\nu^{(p-\mu)} - e_\nu\right)\right) \\
\text{vanishes identically,} \\
\quad \vartheta\left(\overset{p}{\underset{1}{v}}\left(\sum_1^{m+1} \alpha_\nu^{(p+2-\mu)} - \sum_1^m u_\nu^{(p-\mu)} - e_\nu\right)\right) \\
\text{but does not vanish identically.}
\end{array}
\right.
\]
This latter function, regarded as a function of $\zeta_{p+1}$, then vanishes for $\varepsilon_{p-1}$, $\varepsilon_{p-2}$, \dots, $\varepsilon_{p-m}$, and therefore in addition for $p-m$ points; and if one denotes these by $\eta_1$, $\eta_2$, \dots, $\eta_{p-m}$, then
\[
\left(\overset{p}{\underset{1}{v}}\left(-\sum_{p-m+1}^p \alpha_\nu^{(\mu)} + e_\nu\right)\right) \equiv \left(\overset{p}{\underset{1}{v}}\left(\sum_1^{p-m} \alpha_\nu^{(\mu)}\right)\right)
\]
and these points $\eta_1$, $\eta_2$, \dots, $\eta_{p-m}$ can be determined in only one way such that this congruence is satisfied, because otherwise the function would vanish for more than $p$ points. The same function, regarded as a function of $z_{p-1}$, vanishes, besides for $\eta_{p+1}$, $\eta_p$, \dots, $\eta_{p-m+1}$, also for $p-m-1$ points, and if one denotes these by $\varepsilon_1$, $\varepsilon_2$, \dots, $\varepsilon_{p-m-1}$, then
\[
\left(\overset{p}{\underset{1}{v}}\left(-\sum_{p-m}^{p-2} u_\nu^{(\mu)} - e_\nu\right)\right) \equiv \left(\overset{p}{\underset{1}{v}}\left(\sum_1^{p-m-1} u_\nu^{(\mu)}\right)\right),
\]
and the points $\varepsilon_1$, $\varepsilon_2$, \dots, $\varepsilon_{p-m-1}$ are completely determined by this congruence.

Under the assumption made in (1.), therefore, in order to satisfy the congruences
\[
\tag{(2.)}
\left(\overset{p}{\underset{1}{v}}(e_\nu)\right) \equiv \left(\overset{p}{\underset{1}{v}}\left(\sum_1^p \alpha_\nu^{(\mu)}\right)\right)
\]
and
\[
\tag{(3.)}
\left(\overset{p}{\underset{1}{v}}(-e_\nu)\right) \equiv \left(\overset{p}{\underset{1}{v}}\left(\sum_1^{p-2} u_\nu^{(\mu)}\right)\right)
\]
$m$ of the points $\eta$ and $m-1$ of the points $\varepsilon$ can be chosen arbitrarily, but thereby the remainder are determined. Evidently these theorems also hold conversely, i.e.\ the function vanishes if one of these conditions is fulfilled. If therefore congruence (2.) is solvable in more than one way, then congruence (3.) is also solvable, and if of the points $\eta$, $m$ but no more can be chosen arbitrarily, then of the points $\varepsilon$, $m-1$ can be chosen arbitrarily and thereby the remainder are determined, and conversely.

\origpage{166}
In an entirely similar way it follows that if
\[
\vartheta\left(\overset{p}{\underset{1}{v}}(r_\nu)\right) = 0
\]
the congruences
\[
\tag{(4.)}
\left(\overset{p}{\underset{1}{v}}(r_\nu)\right) \equiv \left(\overset{p}{\underset{1}{v}}\left(\sum_1^{p-1} \alpha_\nu^{(\mu)}\right)\right),
\]
\[
\tag{(5.)}
\left(\overset{p}{\underset{1}{v}}(-r_\nu)\right) \equiv \left(\overset{p}{\underset{1}{v}}\left(\sum_1^{p-1} u_\nu^{(\mu)}\right)\right)
\]
are always solvable; and indeed, of the points $\eta$ as well as of the points $\varepsilon$, $m$ can be chosen arbitrarily, and thereby the remaining $p-1-m$ are determined, if
\[
\vartheta\left(\overset{p}{\underset{1}{v}}\left(\sum_1^m u_\nu^{(\mu)} - \sum_1^m \alpha_\nu^{(\mu)} + r_\nu\right)\right)
\]
is identically equal to zero, but
\[
\vartheta\left(\overset{p}{\underset{1}{v}}\left(\sum_1^{m+1} u_\nu^{(\mu)} - \sum_1^{m+1} \alpha_\nu^{(\mu)} + r_\nu\right)\right)
\]
is not identically equal to zero, where the case $m=0$ is not excluded. This theorem can also be reversed. If therefore of the points $\eta$, $m$ and no more can be chosen arbitrarily, its premise is satisfied; and consequently of the points $\varepsilon$ also, $m$ and no more can be chosen arbitrarily.

\section*{4.}

\[
\tag{(1.)}
\left\{
\begin{array}{l}
\text{If we denote the derivative of} \\
\quad \vartheta(v_1, v_2, \dots, v_p) \\
\text{with respect to } v_\nu \text{ by } \vartheta'_\nu\text{, the second derivative with respect to } v_\nu \text{ and } v_\mu \text{ by } \vartheta''_{\nu,\mu}\text{, etc.,}
\end{array}
\right.
\]
then, if $\vartheta\left(\overset{p}{\underset{1}{v}}(u_\nu^{(1)} - \alpha_\nu^{(1)} + r_\nu)\right)$ vanishes identically for every value of $z_1$ and $\zeta_1$, all functions $\vartheta'\left(\overset{p}{\underset{1}{v}}(r_\nu)\right)$ are equal to zero. In fact, the equation
\[
\vartheta\left(\overset{p}{\underset{1}{v}}(u_\nu^{(1)} - \alpha_\nu^{(1)} + r_\nu)\right) = 0,
\]
\origpage{167}
when $s_1$ and $z_1$ differ infinitely little from $\sigma_1$ and $\zeta_1$, passes over into the equation
\[
\sum_1^p \vartheta'_\mu\left(\overset{p}{\underset{1}{v}}(r_\nu)\right) d\alpha_\mu^{(1)} = 0.
\]
If we assume that
\[
du_\mu = \frac{\varphi_\mu(s, z)\,\partial z}{\displaystyle\frac{\partial F}{\partial s}}
\]
then after omitting the factor $\displaystyle\frac{\partial \zeta_1}{\displaystyle\frac{\partial F(\sigma_1, \zeta_1)}{\partial \sigma_1}}$ this equation transforms into
\[
\sum_1^p \vartheta'_\mu\left(\overset{p}{\underset{1}{v}}(r_\nu)\right) \varphi_\mu(\sigma_1, \zeta_1) = 0;
\]
and since between the functions $\varphi$ no linear equation with constant coefficients subsists, it follows from this that all first derivatives of $\vartheta(v_1, v_2, \dots, v_p)$ must vanish for $\overset{p}{\underset{1}{v}}(v_\nu = r_\nu)$.

To prove the converse theorem, we assume that $\overset{p}{\underset{1}{v}}(v_\nu = r_\nu)$ and $\overset{p}{\underset{1}{v}}(v_\nu = t_\nu)$ are two systems of values for which the function $\vartheta$ vanishes without vanishing identically for $\overset{p}{\underset{1}{v}}(v_\nu = u_\nu^{(1)} - \alpha_\nu^{(1)} + r_\nu)$ and $\overset{p}{\underset{1}{v}}(v_\nu = u_\nu^{(1)} - \alpha_\nu^{(1)} + t_\nu)$, and form the expression
\[
\tag{(2.)}
\frac{\vartheta\left(\overset{p}{\underset{1}{v}}(u_\nu^{(1)} - \alpha_\nu^{(1)} + r_\nu)\right)\vartheta\left(\overset{p}{\underset{1}{v}}(\alpha_\nu^{(1)} - u_\nu^{(1)} + r_\nu)\right)}{\vartheta\left(\overset{p}{\underset{1}{v}}(u_\nu^{(1)} - \alpha_\nu^{(1)} + t_\nu)\right)\vartheta\left(\overset{p}{\underset{1}{v}}(\alpha_\nu^{(1)} - u_\nu^{(1)} + t_\nu)\right)}.
\]
If we consider this expression as a function of $z_1$, it follows that it is an algebraic function of $z_1$, and indeed a rational function of $s_1$ and $z_1$, since denominator and numerator are continuous in $T''$ and acquire the same factors at the crosscuts. For $z_1 = \zeta_1$ and $s_1 = \sigma_1$ denominator and numerator become infinitely small of the second order, so that the function remains finite; but the other values for which denominator or numerator vanish are, as proved above, completely determined by the values of the quantities $r$ and of the quantities $t$, and thus entirely independent of $\zeta_1$. Since now an algebraic function is determined up to a
\origpage{168}
constant factor by the values for which it becomes zero and infinite, the expression is equal to a rational function of $s_1$ and $z_1$ independent of $\zeta_1$, $\chi(s_1, z_1)$, multiplied by a constant, i.e.\ a quantity independent of $z_1$. Since the expression is symmetrical with respect to the systems of quantities $(s_1, z_1)$ and $(\sigma_1, \zeta_1)$, this constant is equal to $\chi(\sigma_1, \zeta_1)$, multiplied by a quantity $A$ that is also independent of $\zeta_1$. If one now sets $\sqrt{A}\cdot\chi(s, z) = \varrho(s, z)$, one obtains for our expression (2.) the value
\[
\tag{(3.)}
\varrho(s_1, z_1)\,\varrho(\sigma_1, \zeta_1)
\]
where $\varrho(s, z)$ is a rational function of $s$ and $z$.

In order to determine this, one needs only to let $\zeta_1 = z_1$ and $\sigma_1 = s_1$; there then results
\[
(\varrho(s_1, z_1))^2 = \left\{\frac{\sum_\mu \vartheta'_\mu\left(\overset{p}{\underset{1}{v}}(r_\nu)\right) du_\mu^{(1)}}{\sum_\mu \vartheta'_\mu\left(\overset{p}{\underset{1}{v}}(t_\nu)\right) du_\mu^{(1)}}\right\}^2
\]
or, after extracting the square root and canceling the factor $\displaystyle\frac{dz_1}{\displaystyle\frac{\partial F(s_1, z_1)}{\partial s_1}}$,
\[
\tag{(4.)}
\varrho(s_1, z_1) = \pm \frac{\sum_\mu \vartheta'_\mu\left(\overset{p}{\underset{1}{v}}(r_\nu)\right) \varphi_\mu(s_1, z_1)}{\sum_\mu \vartheta'_\mu\left(\overset{p}{\underset{1}{v}}(t_\nu)\right) \varphi_\mu(s_1, z_1)}.
\]
One therefore has from (3.) and (4.) the equation
\[
\tag{(5.)}
\left\{
\begin{aligned}
& \frac{\vartheta\left(\overset{p}{\underset{1}{v}}(u_\nu^{(1)} - \alpha_\nu^{(1)} + r_\nu)\right)\vartheta\left(\overset{p}{\underset{1}{v}}(\alpha_\nu^{(1)} - u_\nu^{(1)} + r_\nu)\right)}{\vartheta\left(\overset{p}{\underset{1}{v}}(u_\nu^{(1)} - \alpha_\nu^{(1)} + t_\nu)\right)\vartheta\left(\overset{p}{\underset{1}{v}}(\alpha_\nu^{(1)} - u_\nu^{(1)} + t_\nu)\right)} \\
&= \frac{\sum_\mu \vartheta'_\mu\left(\overset{p}{\underset{1}{v}}(r_\nu)\right) \varphi_\mu(s_1, z_1)}{\sum_\mu \vartheta'_\mu\left(\overset{p}{\underset{1}{v}}(t_\nu)\right) \varphi_\mu(s_1, z_1)} \cdot \frac{\sum_\mu \vartheta'_\mu\left(\overset{p}{\underset{1}{v}}(r_\nu)\right) \varphi_\mu(\sigma_1, \zeta_1)}{\sum_\mu \vartheta'_\mu\left(\overset{p}{\underset{1}{v}}(t_\nu)\right) \varphi_\mu(\sigma_1, \zeta_1)}.
\end{aligned}
\right.
\]
From this equation it follows that $\vartheta\left(\overset{p}{\underset{1}{v}}(u_\nu^{(1)} - \alpha_\nu^{(1)} + r_\nu)\right)$ must be equal to zero for every value
\origpage{169}
of $z_1$ and $\zeta_1$ if the first derivatives of the function $\vartheta(v_1, v_2, \dots, v_p)$ all vanish for $\overset{p}{\underset{1}{v}}(v_\nu = r_\nu)$.

\section*{5.}

If
\[
\tag{(1.)}
\vartheta\left(\overset{p}{\underset{1}{v}}\left(\sum_1^m \alpha_\nu^{(\mu)} - \sum_1^m u_\nu^{(\mu)} + r_\nu\right)\right)
\]
vanishes identically, i.e.\ for all values of $\overset{m}{\underset{1}{\mu}}(\sigma_\mu, \zeta_\mu)$ and $\overset{m}{\underset{1}{\mu}}(s_\mu, z_\mu)$, then by the route indicated above one finds, first, by letting $\zeta_m = z_m$, $\sigma_m = s_m$, that the first derivatives of the function $\vartheta(v_1, v_2, \dots, v_p)$ all vanish for $\overset{p}{\underset{1}{v}}\left(v_\nu = \displaystyle\sum_1^{m-1} \alpha_\nu^{(\mu)} - \displaystyle\sum_1^{m-1} u_\nu^{(\mu)} + r_\nu\right)$; then, by letting $\zeta_{m-1} - z_{m-1}$, $\sigma_{m-1} - s_{m-1}$ become infinitely small, that for $\overset{p}{\underset{1}{v}}\left(v_\nu = \displaystyle\sum_1^{m-2} \alpha_\nu^{(\mu)} - \displaystyle\sum_1^{m-2} u_\nu^{(\mu)} + r_\nu\right)$ the second derivatives also all vanish; and evidently it follows in general that the derivatives of the $n^{\text{th}}$ order all vanish for $\overset{p}{\underset{1}{v}}\left(v_\nu = \displaystyle\sum_1^{m-n} \alpha_\nu^{(\mu)} - \displaystyle\sum_1^{m-n} u_\nu^{(\mu)} + r_\nu\right)$, whatever values the quantities $z$ and the quantities $\zeta$ may have.

It follows from this that under the present assumption (1.), for $\overset{p}{\underset{1}{v}}(v_\nu = r_\nu)$ the first through $m^{\text{th}}$ derivatives of the function $\vartheta(v_1, v_2, \dots, v_p)$ are all equal to zero.

To show that this theorem also holds conversely, we first prove that if $\vartheta\left(\overset{p}{\underset{1}{v}}\left(\displaystyle\sum_1^{m-1} \alpha_\nu^{(\mu)} - \displaystyle\sum_1^{m-1} u_\nu^{(\mu)} + r_\nu\right)\right)$ vanishes identically and the quantities $\vartheta^{(m)}\left(\overset{p}{\underset{1}{v}}(r_\nu)\right)$ are all equal to zero, then $\vartheta\left(\overset{p}{\underset{1}{v}}\left(\displaystyle\sum_1^m \alpha_\nu^{(\mu)} - \displaystyle\sum_1^m u_\nu^{(\mu)} + r_\nu\right)\right)$ must also vanish identically, and for this purpose we generalize the equation (\S.~4, (5.)).

We assume that $\vartheta\left(\overset{p}{\underset{1}{v}}\left(\displaystyle\sum_1^{m-1} u_\nu^{(\mu)} - \displaystyle\sum_1^{m-1} \alpha_\nu^{(\mu)} + r_\nu\right)\right)$ vanishes identically, but $\vartheta\left(\overset{p}{\underset{1}{v}}\left(\displaystyle\sum_1^m u_\nu^{(\mu)} - \displaystyle\sum_1^m \alpha_\nu^{(\mu)} + r_\nu\right)\right)$ does not vanish identically, retain with respect
\origpage{170}
to the quantities $t$ the earlier assumption, and consider the expression
\[
\tag{(2.)}
\frac{\vartheta\left(\overset{p}{\underset{1}{v}}\left(\sum_1^m u_\nu^{(\mu)} - \sum_1^m \alpha_\nu^{(\mu)} + r_\nu\right)\right) \vartheta\left(\overset{p}{\underset{1}{v}}\left(\sum_1^m \alpha_\nu^{(\mu)} - \sum_1^m u_\nu^{(\mu)} + r_\nu\right)\right) \prod_{\varrho\varrho'} \vartheta\left(\overset{p}{\underset{1}{v}}(u_\nu^\varrho - u_\nu^{\varrho'} + t_\nu)\right) \vartheta\left(\overset{p}{\underset{1}{v}}(\alpha_\nu^\varrho - \alpha_\nu^{\varrho'} + t_\nu)\right)}{\left(\overset{m}{\underset{1}{\Pi}}\right)^2 \vartheta\left(\overset{p}{\underset{1}{v}}(u_\nu^\varrho - \alpha_\nu^{\varrho'} + t_\nu)\right) \vartheta\left(\overset{p}{\underset{1}{v}}(\alpha_\nu^\varrho - u_\nu^{\varrho'} + t_\nu)\right)}
\]
In this expression, under the product signs all values from 1 to $m$ are to be taken for $\varrho$ as well as for $\varrho'$, but in the numerator the cases where $\varrho = \varrho'$ would occur are to be omitted.

If we consider this expression as a function of $z_1$, it follows that it acquires the factor 1 at the crosscuts and is consequently an algebraic function of $z_1$. For $z_1 = \zeta_\varrho$ and $s_1 = \sigma_\varrho$ denominator and numerator become infinitely small of the second order, the fraction therefore remains finite; but the other values for which numerator and denominator vanish are, as proved above (\S.~3), completely determined by the quantities $\overset{m}{\underset{2}{\mu}}(s_\mu, z_\mu)$, the quantities $r$, and the quantities $t$, and consequently entirely independent of the quantities $\zeta$. Since the expression is now a symmetrical function of the quantities $z$, the same holds for any arbitrary $z_\mu$: it is an algebraic function of $z_\mu$, and the values of this quantity for which it becomes infinitely large or infinitely small are independent of the quantities $\zeta$. It is therefore equal to an algebraic function of the quantities $z$ independent of the quantities $\zeta$, $\chi(z_1, z_2, \dots, z_m)$, multiplied by a factor independent of the quantities $z$. But since it remains unchanged when one interchanges the quantities $z$ with the quantities $\zeta$, this factor is equal to $\chi(\zeta_1, \zeta_2, \dots, \zeta_m)$, multiplied by a constant $A$ independent of the quantities $z$ and the quantities $\zeta$; and we can therefore, if we set $\sqrt{A}\cdot\chi(z_1, z_2, \dots, z_m) = \psi(z_1, z_2, \dots, z_m)$, give to our expression (2.) the form
\[
\tag{(3.)}
\psi(z_1, z_2, \dots, z_m)\,\psi(\zeta_1, \zeta_2, \dots, \zeta_m)
\]
where $\psi(z_1, z_2, \dots, z_m)$ is an algebraic function of the quantities $z$ independent of the quantities $\zeta$, which in consequence of its mode of branching must be expressible rationally in $\overset{m}{\underset{1}{\mu}}(s_\mu, z_\mu)$. If now one lets the points $\eta$ coincide with the points $\varepsilon$, so that the quantities $\zeta_\mu - z_\mu$ and the quantities $\sigma_\mu - s_\mu$ all become infinitely small, there results, if one denotes the derivatives
\origpage{171}
of $\vartheta(v_1, v_2, \dots, v_p)$ as above (\S.~4, (1.)),
\[
\tag{(4.)}
\psi(z_1, z_2, \dots, z_m) = \pm \frac{\left(\overset{p}{\underset{1}{\Sigma}}\right)^m \vartheta^{(m)}_{\nu_1, \nu_2, \dots, \nu_m}\left(\overset{p}{\underset{1}{\varrho}}(r_\varrho)\right) du_{\nu_1}^{(1)} du_{\nu_2}^{(2)} \dots du_{\nu_m}^{(m)}}{\displaystyle\prod_{\mu=1}^{\mu=m} \prod_{\nu=1}^{\nu=p} \vartheta'_\nu\left(\overset{p}{\underset{1}{\varrho}}(t_\varrho)\right) du_\nu^{(\mu)}},
\]\ednote{The second product sign in the denominator is printed as $\Pi$; from the context, a summation over $\nu$ is required. Reproduced here as printed.}
where the summations in the numerator refer to $\nu_1$, $\nu_2$, \dots, $\nu_m$. It is hardly necessary to remark that the choice of sign is indifferent, since it has no influence upon the value of $\psi(z_1, z_2, \dots, z_m)\,\psi(\zeta_1, \zeta_2, \dots, \zeta_m)$, and that instead of the quantities $du_1^{(\mu)}$, $du_2^{(\mu)}$, \dots, $du_p^{(\mu)}$ one can also introduce, simultaneously in numerator and denominator, the quantities $\varphi_1(s_\mu, z_\mu)$, $\varphi_2(s_\mu, z_\mu)$, \dots, $\varphi_p(s_\mu, z_\mu)$ proportional to them.

From the equation contained in (2.), (3.), and (4.), which is proved for the case where
\[
\vartheta\left(\overset{p}{\underset{1}{v}}\left(\sum_1^{m-1} u_\nu^{(\mu)} - \sum_1^{m-1} \alpha_\nu^{(\mu)} + r_\nu\right)\right)
\]
is equal to zero and
\[
\vartheta\left(\overset{p}{\underset{1}{v}}\left(\sum_1^m u_\nu^{(\mu)} - \sum_1^m \alpha_\nu^{(\mu)} + r_\nu\right)\right)
\]
is different from zero, it follows that
\[
\vartheta\left(\overset{p}{\underset{1}{v}}\left(\sum_1^m u_\nu^{(\mu)} - \sum_1^m \alpha_\nu^{(\mu)} + r_\nu\right)\right)
\]
cannot be different from zero if the functions $\vartheta^{(m)}\left(\overset{p}{\underset{1}{v}}(r_\nu)\right)$ are all equal to zero.

If therefore the functions $\vartheta^{(m+1)}\left(\overset{p}{\underset{1}{v}}(r_\nu)\right)$ are all equal to zero, then from the validity of the equation
\[
\vartheta\left(\overset{p}{\underset{1}{v}}\left(\sum_1^n u_\nu^{(\mu)} - \sum_1^n \alpha_\nu^{(\mu)} + r_\nu\right)\right) = 0
\]
for $n=m$ follows its validity for $n=m+1$. If therefore the equation holds for $n=0$, or $\vartheta\left(\overset{p}{\underset{1}{v}}(r_\nu)\right) = 0$, and the first through $m^{\text{th}}$ derivatives of the function $\vartheta\left(\overset{p}{\underset{1}{v}}(v_\nu)\right)$ all vanish for $\overset{p}{\underset{1}{v}}(v_\nu = r_\nu)$, but the $(m+1)^{\text{th}}$ do not all vanish,
\origpage{172}
then the equation also holds for all greater values of $n$ up to $n=m$, but not for $n=m+1$; for from $\displaystyle\vartheta\left(\overset{p}{\underset{1}{v}}\left(\sum_1^{m+1} u_\nu^{(\mu)} - \sum_1^{m+1} \alpha_\nu^{(\mu)} + r_\nu\right)\right) = 0$ it would follow, as we had already found previously, that the quantities $\vartheta^{(m+1)}\left(\overset{p}{\underset{1}{v}}(r_\nu)\right)$ would all have to vanish.

\section*{6.}

If we combine what has just been proved with what went before, we obtain the following result:

If $\vartheta(r_1, r_2, \dots, r_p) = 0$, then $(p-1)$ points $\eta_1$, $\eta_2$, \dots, $\eta_{p-1}$ can be determined such that
\[
(r_1, r_2, \dots, r_p) \equiv \left(\sum_1^{p-1} \alpha_1^{(\mu)}, \sum_1^{p-1} \alpha_2^{(\mu)}, \dots, \sum_1^{p-1} \alpha_p^{(\mu)}\right);
\]
and conversely.

If, besides the function $\vartheta(v_1, v_2, \dots, v_p)$, its first through $m^{\text{th}}$ derivatives are also all equal to zero for $v_1 = r_1$, $v_2 = r_2$, \dots, $v_p = r_p$, but the $(m+1)^{\text{th}}$ are not all equal to zero, then $m$ of these points $\eta$ can be chosen arbitrarily without changing the quantities $r$, and thereby the remaining $p-1-m$ are completely determined.

And conversely:

If $m$ and no more of the points $\eta$ can be chosen arbitrarily without changing the quantities $r$, then, besides the function $\vartheta(v_1, v_2, \dots, v_p)$, its first through $m^{\text{th}}$ derivatives are also all equal to zero for $v_1 = r_1$, $v_2 = r_2$, \dots, $v_p = r_p$, but the $(m+1)^{\text{th}}$ are not all equal to zero.

The complete investigation of all special cases that can occur in the vanishing of a $\vartheta$-function was less necessary on account of the particular systems of equally branched algebraic functions for which these cases arise, than rather because without this investigation gaps would arise in the proof of the theorems founded upon our theorem on the vanishing of a $\vartheta$-function.

Göttingen, in October 1865.

\end{document}
